\emph{(i)} Off the bifurcation surface write $\mathcal{H} = \{S = 0\}$ locally, with physical gradient $n = DS \ne 0$; since $K$ is normal to $\mathcal{H}$ and nonzero there, $K = \phi\,n$ on $\mathcal{H}$ with $\phi \ne 0$. For physical vectors $X, Y$ tangent to $\mathcal{H}$ (so $Y\cdot n = 0$, and the derivative of $K$ along $X$ only involves $K$ on $\mathcal{H}$),
\[
Y\cdot\bigl((X\cdot D)K\bigr) = \phi\,Y\cdot\bigl((X\cdot D)n\bigr) + (X\cdot D\phi)\,(Y\cdot n) = \phi\,Y\cdot\bigl((X\cdot D)DS\bigr),
\]
which is symmetric in $X, Y$ because $\Omega$ is torsion-free (the covariant Hessian of a scalar is symmetric). By Fact \claimref{synthesis.imported-gtg-dictionary}(vi) the same quantity is $Y\cdot(X\cdot F_{K})$, antisymmetric in $X, Y$. Hence $Y\cdot(X\cdot F_{K}) = 0$ for all tangent $X, Y$. The tangent space of $\mathcal{H}$ is $k^{\perp} = \mathrm{span}\{k, a, b\}$ (it contains $k$, which is null). Expand $F_{K}$ in the basis $k\wedge m$, $k\wedge a$, $k\wedge b$, $m\wedge a$, $m\wedge b$, $a\wedge b$: with $X\cdot(u\wedge v) = (X\cdot u)v - (X\cdot v)u$ one finds that $k\wedge m$, $k\wedge a$, $k\wedge b$ satisfy the condition (for $X \in k^{\perp}$, $X\cdot(k\wedge v) = -(X\cdot v)k$, a multiple of $k$, orthogonal to every $Y \in k^{\perp}$), while the pairs $(X, Y) = (a, b)$, $(b, a)$ and $(a, k)$, $(b, k)$ detect the coefficients of $a\wedge b$, $m\wedge a$ and $m\wedge b$ respectively ($k\cdot(a\cdot(m\wedge a)) = -Q(a)\,\ip{k}{m} \ne 0$, $b\cdot(a\cdot(a\wedge b)) = Q(a)Q(b) \ne 0$, and so on). So the solution space is exactly $\mathrm{span}\{k\wedge m, k\wedge a, k\wedge b\}$; the check confirms the rank $3$ of the nine-by-six constraint matrix and the independence of the three bivectors exactly, for $k = e_{0} + e_{1}$ and for the tilted null vector $k = 5e_{0} + 3e_{1} + 4e_{2}$, and that $a\wedge b$, $m\wedge a$, $m\wedge b$ violate the condition. Write $F_{K} = c\,k\wedge m + k\wedge\omega$ with $\omega \in \mathrm{span}\{a, b\}$. Since $k\cdot(k\wedge m) = Q(k)m - \ip{k}{m}k = -k$ and $k\cdot(k\wedge\omega) = Q(k)\omega - \ip{k}{\omega}k = 0$, the surface-gravity condition $(K\cdot D)K = K\cdot F_{K} = \kappa K$ with $K = \phi k$ gives $-c\phi k = \kappa\phi k$, so $c = -\kappa$. Finally $(k\wedge m)^{2} = \ip{k}{m}^{2} - Q(k)Q(m) = 1$; with $k\wedge m = km - 1$ one has $(k\wedge m)k\omega = k\omega$ and $k\omega(k\wedge m) = -k\omega$ (since $k\omega k = -kk\omega = 0$), so the two anticommute; and $(k\omega)^{2} = -kk\omega\omega = 0$. Hence $F_{K}^{2} = \kappa^{2}$. All identities are checked exactly for $\kappa = \tfrac14, -\tfrac27, 3$, two rotation one-forms and both null frames.

\emph{(ii)} Choose near the generator a frame $\{d_{a}\}$ dragged by the Killing flow, $[K, d_{a}] = 0$; it stays orthonormal since $K$ is Killing. For a dragged vector field $A$, torsion-freeness gives $(K\cdot D)A = (A\cdot D)K = A\cdot F_{K} = -F_{K}\times A$, and the components of $F_{K}$ in the dragged frame are constant along the orbit (the Killing flow preserves $D\wedge K$). A vector $P$ parallel along the orbit, $(K\cdot D)P = 0$, has dragged components obeying $\dot P = F_{K}\times P$, so $P(t) = S_{t}P(0)S_{t}^{\sim}$ with $S_{t} = \exp(tF_{K}/2)$: parallel frames are the dragged frames turned by $S_{t}$, and since $S_{t}$ commutes with $F_{K}$ the components of $F_{K}$ are constant in the parallel frame as well ($F_{K}$ is parallel along the orbit). A spinor field invariant under the flow has constant dragged components; relative to the parallel spin frame (Fact \claimref{synthesis.imported-gtg-dictionary}(ii), $dU/dt = -\tfrac12\Omega U$) its values are $\psi(t) = S_{t}^{\sim}\psi(0) = R_{t}\psi(0)$, $R_{t} = \exp(-tF_{K}/2)$. As $F_{K}^{2} = \kappa^{2}$, $R_{t} = \cosh\tfrac{\kappa t}{2} - \sinh\tfrac{\kappa t}{2}F_{K}/\kappa$, and $R_{t}R_{s} = R_{t+s}$, $R_{t}R_{t}^{\sim} = 1$. With $\Lambda(a) = a\gamma_{0}/(2a_{0})$ of Proposition \claimref{field-theory.second-quantisation}(1), writing $\gamma_{0} = \alpha k + \beta m$ (it lies in the plane of $k$ and $m$), $\Lambda(m)k = m\gamma_{0}k/(2m_{0}) = (\alpha\,mkk + \beta\,mmk)/(2m_{0}) = 0$; with $k\wedge m = km - 1$ this gives $\Lambda(m)(k\wedge m) = -\Lambda(m)$, $\Lambda(m)k\omega = 0$, hence $\Lambda(m)F_{K} = \kappa\Lambda(m)$ and $\Lambda(m)R_{t} = (\cosh\tfrac{\kappa t}{2} - \sinh\tfrac{\kappa t}{2})\Lambda(m) = e^{-\kappa t/2}\Lambda(m)$. Likewise $kF_{K} = -\kappa k(km - 1) + kk\omega = \kappa k$ and $kR_{t} = e^{-\kappa t/2}k$; a null rotation $1 + \tfrac12 ks$ about $k$ (an element of the abelian group of null rotations fixing $k$, Proposition \claimref{transport.group-null-rotations}; it changes $m$) leaves $k$ fixed from the left, so the data $k\psi$ do not depend on the choice of $m$. On the vectors, $R_{t}kR_{t}^{\sim} = e^{\kappa t}k$: relative to parallel frames the invariant vector $K$ grows as $e^{\kappa t}$, which is the statement $K = \kappa u\,k_{\mathrm{aff}}$ with $u \mapsto e^{\kappa t}u$ of the boost $\beta_{t}$ of the setup of Theorem \claimref{synthesis.positive-null-energy-gives-r-premises}. The null-rotation part maps the canonical range into the constrained one: $\Lambda(m)$ is a left eigen-element of $F_{K}$ ($\Lambda(m)F_{K} = \kappa\Lambda(m)$), so the canonical component of $R_{t}\psi$ is $e^{-\kappa t/2}\Lambda(m)\psi$ whatever $\Lambda(k)\psi$ is, while $F_{K}\Lambda(m) \ne \kappa\Lambda(m)$ when $\omega \ne 0$. A transverse rotation $\exp(\theta\,a\wedge b/2)$ would commute with $\Lambda(m)$ without being a scalar on its range (a phase on the two chirality channels, checked as a control); (i) excludes it. The check verifies all of this exactly at $e^{\kappa t/2} = \tfrac32, \tfrac25, \tfrac73$, for boosts composed with null rotations in either order, and on the closed form.

\emph{(iii)} Theorem \claimref{synthesis.boost-evolution-on-generators-is-one-unitary-group}(ii): a value factor $e^{-\kappa t/2}$ rescales the canonical kernel by $e^{-\kappa t}$, and the transport $u \mapsto e^{\kappa t}u$ rescales the delta kernel by $e^{\kappa t}$, so the composite preserves it and is the dilation $D_{\kappa t}$ of Proposition \claimref{synthesis.affine-cell-boost-spreads-without-scrambling}; for $\kappa < 0$ (the cosmological orientation, $u \mapsto e^{-|\kappa|t}u$) the same holds with $t \mapsto -t$. The check verifies the kernel identity on affine cells at three rapidities, with each factor alone failing. Since (ii) shows the value factor is exactly the boost $B_{\kappa t}$ on $\Lambda(m)$, with no phase, the identification argued in that theorem's (vi) is derived.

\emph{(iv)} For $\hb(a) = a - \beta(a\cdot e_{t})e_{r}$ the derived metric (Definition \claimref{obj.synthesis.gtg-flat-spacetime-gauge-fields}(d)) has the orthonormal coframe $\theta^{0} = dt$, $\theta^{1} = dr + \beta\,dt$, $\theta^{2} = r\,d\vartheta$, $\theta^{3} = r\sin\vartheta\,d\varphi$, with $Q(\beta) = 1 - f$. Its torsion-free metric connection, the rotation-gauge field read as connection forms (Fact \claimref{synthesis.imported-gtg-dictionary}(ii)), is $\omega^{1}{}_{0} = -\beta'\theta^{1}$, $\omega^{2}{}_{0} = -(\beta/r)\theta^{2}$, $\omega^{3}{}_{0} = -(\beta/r)\theta^{3}$, $\omega^{2}{}_{1} = \theta^{2}/r$, $\omega^{3}{}_{1} = \theta^{3}/r$, $\omega^{3}{}_{2} = (\cot\vartheta/r)\theta^{3}$: by hand from $d\theta^{1} = \beta'\theta^{1}\wedge\theta^{0}$, $d\theta^{2} = (\theta^{1} - \beta\theta^{0})\wedge\theta^{2}/r$, $d\theta^{3} = (\theta^{1} - \beta\theta^{0})\wedge\theta^{3}/r + (\cot\vartheta/r)\theta^{2}\wedge\theta^{3}$, and confirmed in the check by Cartan's first structure equation $d\theta^{a} = -\omega^{a}{}_{b}\wedge\theta^{b}$ and metric compatibility at three sample points with independent rational $\beta, \beta', r, \cot\vartheta$ (a connection with a perturbed radial coefficient fails). Contracting with $K = \partial_{t} = E_{0} + \beta E_{1}$ kills every form except $\omega^{1}{}_{0}$, which gives $-\beta\beta' = f'/2$: $\Omega(e_{t}) = \tfrac{f'}{2}\sigma_{r}$ with no angular part, at every radius, as Proposition \claimref{singularity.plane-rate} found for Schwarzschild. The frame is $t$-independent, so it is dragged by $K$ and $\Omega(e_{t}) = -F_{K}$ by (ii). At a horizon, $f = 0$ and $Q(\beta) = 1$; the radial null rays have $dr/dt = \pm1 - \beta$, so for $\beta = 1$ the $t$-line is the outgoing generator, $K = e_{t} + e_{r} = k$ is null, $m = \tfrac12(e_{t} - e_{r})$, $k\wedge m = e_{r}e_{t} = \sigma_{r}$, and $(K\cdot D)K = \Omega(e_{t})\times K = \tfrac{f'}{2}K$, so $\kappa = f'(r_{h})/2$ and $\omega = 0$. Instances checked exactly: Schwarzschild $\beta = \sqrt{2M/r}$, $\beta(2M) = 1$, $\beta'(2M) = -1/4M$, $\kappa = 1/4M$ for $M = 1, \tfrac32, \tfrac17$; de Sitter $\beta = Hr$, $\kappa = -H$ at $r = 1/H$ for $H = 1, \tfrac23$; Schwarzschild-de Sitter $f = 1 - 2M/r - H^{2}r^{2}$ with $H^{2} = \tfrac17$, $M = \tfrac37$, whose horizons are $r = 1$ and $r = 2$ ($f$ vanishes there exactly), with $\beta' = -f'/2$ at $\beta = 1$ and $\kappa = \tfrac27$, $-\tfrac{5}{28}$. Rotating horizons are covered by (i) to (iii) without a gauge computation (the same block-triangular mechanism applies to any null rotation about the generator, whose canonical block is the identity): their Killing bivector has $\omega \ne 0$ off the bifurcation surface, which (ii) shows does not reach the canonical component.

Scope: the statement is local on the horizon off the bifurcation surface (where $K = 0$ and the generator parameter degenerates; there $F_{K}$ is $-\kappa$ times the binormal). It assumes a Killing horizon with $\kappa \ne 0$ and excludes degenerate (extremal) horizons; it says nothing about horizons without a Killing field (a collapsing hole's horizon), for which the stored local-horizon results apply instead.
